\documentclass{article}
\usepackage[utf8]{inputenc}
\usepackage[T1]{fontenc}
\usepackage{geometry}
\usepackage{graphicx}
\usepackage{float}
\usepackage{hyperref}
\usepackage[english,indonesian]{babel}

\title{Optimasi Pemrosesan Template LaTeX Menggunakan Pendekatan Berbasis AI}
\author{Academic Support Agent}
\date{\today}

\begin{document}
\maketitle

\begin{abstract}

\end{abstract}

\chapter{PENDAHULUAN}


\chapter{TINJAUAN PUSTAKA}


% Original content (commented out):
% \input{preambule}
\begin{document}

% UZUPEŁNIAMY PLIK CHAPTER1/CHAPTER.TEX!

\pagestyle{fancy}
\fancyhead[C]{\normalfont{\footnotesize{\textsc{Wydział Elektroniki}}}}
\renewcommand{\headrulewidth}{.5pt}
	\include{chapter1/chapter}
	
\end{document}


\end{document}